{\subsection{Gram and orthogonal matrices}\label{subsec:gram-ortho}
Recall that a real symmetric matrix $M \in \R^{n \times n}$ is positive semidefinite if $\bv^{T} M \bv \ge 0$ for all $\bv \in \R^n$. Since $M$ is real and symmetric, we know that $M$ has a full eigenspace: there exists $\lambda_1, \hdots, \lambda_n \in \R$ and unit vectors $\bv_1, \hdots, \bv_n \in \R^n$ such that 
\[
M = \sum_{i=1}^n \lambda_i \bv \bv_i^{T}.
\]
Since $M$ is positive semidefinite, we know that each $\lambda_i \ge 0$. That is, if we let $V \in \R^{n \times n}$ be the matrix such that the $i$th row is $\sqrt{\lambda_i} \bv_i$, then we have that $M = V^{T} V$. In other words, $M$ is also a \emph{Gram matrix}. It is also straightforward that any Gram matrix is positive semidefinite since if $M = V^{T} V$, then for any $\bv \in \R^n$, we have that $\bv^{T} M \bv = \|V \bv\|_2^2 \ge 0$. We call $V$ a \emph{realization} of the Gram matrix, and we call the columns of $V$ the vectors of the realization.

A crucial fact about realizations is that they generate the same Gram matrix if and only if they are related by an \emph{orthogonal matrix}, a matrix $Q \in \R^{n \times n}$ for which $Q^{T} Q = Q Q^{T} = I_n$. We prove a generalization of this equivalence where the number of vectors differs from the ambient dimension.

\begin{proposition}\label{prop:gram-ortho}
Let $n, m$ be positive integers and let $\bu_1, \hdots, \bu_n, \bv_1, \hdots, \bv_n \in \R^m$ be vectors. The following are equivalent.
\begin{itemize}
\item [(1)] For all $i,j \in [n]$, we have that $\langle \bu_i, \bu_j\rangle = \langle \bv_i, \bv_j \rangle$.
\item [(2)] There exists an orthogonal matrix $Q \in \R^{m \times m}$ such that for all $i \in [n]$, $\bv_i = Q \bu_i$
\end{itemize}
\end{proposition}
\begin{proof}
Proving that (2) implies (1) is straightforward for if $\bv_i = Q\bu_i$ for all $i \in [n]$, we have that
\[
    \langle \bv_i, \bv_j\rangle = \langle Q\bu_i, Q\bu_j\rangle = (Q\bu_i)^{T} Q\bu_j = \bu_i^{T} Q^T Q \bu_j = \langle \bu_i, \bu_j\rangle,
\]
where we use that $Q$ is an orthogonal matrix.

We now prove that (1) implies (2). Let $U := \operatorname{span}\{\bu_i : i \in [n]\} \subset \R^{m}$. Using the Gram-Schmidt algorithm, we can construct an orthonormal basis $\ba_1, \hdots, \ba_m$ of $\R^n$ such that $\ba_1, \hdots, \ba_{\dim U}$ is a basis of $U$. (Assume $\dim U \ge 1$ or else every vector equals $0$ for which we can pick $Q$ to be the identity matrix.)

Also pick $S \subseteq [n]$ such that $\{\bu_{s} : s \in S\}$ is a basis of $U$. There exists a unique linear map $L : U \to \R^m$ for which $L(\bu_{s}) = \bv_s$ for all $s \in S$. We claim the following properties about $L$,
\begin{itemize}
\item[(1)] $L(\bu_i) = \bv_i$ for all $i \in [n]$, so the image of $L$ is $V:= \operatorname{span}\{\bv_i : i \in [n]\}$.
\item[(2)] $L(\ba_1), \hdots, L(\ba_{\dim U})$ is an orthonormal basis of $V$.
\end{itemize}
To verify (1), for any $i \in [n]$, let $\lambda : S \to \R$ be such that $\bu_i = \sum_{s \in S} \lambda(s)\bu_s$. Then, $L(\bu_i) = \sum_{s \in S} \lambda(s)\bv_s$. Note then that
\begin{align*}
\langle L(\bu_i), \bv_i\rangle &= \sum_{s \in S} \lambda(s) \langle \bv_s, \bv_i\rangle = \sum_{s\in S} \lambda(s) \langle \bu_s, \bu_i\rangle = \langle \bu_i, \bu_i\rangle = \langle \bv_i, \bv_i\rangle, \text{ and}\\
\langle L(\bu_i), L(\bu_i)\rangle &= \sum_{s,s' \in S} \lambda(s)\lambda(s') \langle \bv_s, \bv_{s'}\rangle = \lambda(s)\lambda(s') \langle \bu_s, \bu_{s'}\rangle = \langle \bu_i, \bu_i\rangle = \langle \bv_i, \bv_i\rangle.
\end{align*}
This can only happen if $L(\bu_i) = \bv_i$. To verify (2), for each $i \in [\dim U]$, let $\kappa_i : S \to \R$ be such that $\ba_i = \sum_{s\in S} \kappa_i(s) \bu_s$. Then, note that for any $i,j \in [\dim U]$ we have that
\[
\langle \ba_i, \ba_j \rangle = \sum_{s,s'\in S}\kappa_i(s)\kappa_i(s') \langle \bu_s, \bu_{s'}\rangle = \sum_{s,s'\in S}\kappa_i(s)\kappa_i(s') \langle L(\bu_s), L(\bu_{s'})\rangle = \langle L(\ba_i), L(\ba_j), \rangle,
\]
so $L(\ba_1), \hdots, L(\ba_{\dim U})$ is an orthonormal basis of $V$.

To finish, we can extend $L(\ba_1), \hdots, L(\ba_{\dim U})$ into an orthonormal basis $\bb_1, \hdots, \bb_m$ of $\R^m$ such that $\bb_i = L(\ba_i)$ for all $i \in [\dim U]$. We can thus construct a linear map $Q : \R^m \to \R^m$ such that $Q(\ba_i) = \bb_i$ for all $i \in [m]$. The matrix representation of $Q$ is an orthogonal matrix and it sends $\bu_i$ to $\bv_i$ for all $i \in [n]$ because $Q$ extends $L$.
\end{proof}

As an immediate corollary, we recover the equivalence of realizations.
\begin{corollary}\label{cor:gram-ortho}
Let $M = U^{T} U$ and $N = V^{T} V$ be $n \times n$ real matrices. We have that $M= N$ if and only if there exists an orthogonal matrix $Q$ such that $V = QU$. 
\end{corollary}
\begin{proof}
Apply Proposition~\ref{prop:gram-ortho} with $m = n$, $\bu_1, \hdots, \bu_n$ as the columns of $U$, and $\bv_1, \hdots, \bv_n$ as the columns of $V$.
\end{proof}


}
